% section: 漸化式を解く前に

\section{漸化式を解く前に}\label{sec:before}
  漸化式を解く前に,根底にある発想について触れておこう.
  多くの漸化式は係数などを無視して書けば$a_{n+1}=a_n$のような形をしている.
  ここに係数や数式が足されることでどんどん問題は複雑化していく.
  しかしながら不変の事実として,漸化式はある数列の2項の関係を示している.
  だから,\,2行の間で"きれい"な関係が成り立っていればその漸化式は解けるのである.
  この"きれい"というのは,演習を重ねるにつれて理解される,非言語の数学的直感によるものであるが,
  あえて言語化すればその多くは
  \begin{itemize}
    \item 式が基本4パターン（特に等比型）の見た目
    \item 式の左辺が$n+1$に関する式,右辺が$n$に関する式になっている
    \item 式の左辺と右辺が同じ見た目
    \item 特定の（既知の）式になっている
  \end{itemize}
  などである.
  こういった式を作ることを目標に変形をしているということを頭の片隅に置いておくといいだろう.
  この発想は応用パターンの節で頻繁に用いられることになる.

  もちろん,一般項を予測し帰納法によって証明する方法も有効である.
  専らそれは,実験的に最初の何項かを求めている時に発見される.
  あくまで帰納法は漸化式の一般項を得るという状況であれば最終手段である.
  しかし,できないとやらないは大きな違いであるし,
  帰納法を使って証明するかどうかに関わらず,悩んだら実験してみるというのはいかなる問題の時でも常に頭の中に入れておくべきことである.
